\documentclass[11pt,a4paper]{article}

\usepackage[utf8]{inputenc}
\usepackage[T1]{fontenc}
\usepackage[spanish,es-tabla]{babel}
\usepackage{lmodern}
\usepackage[a4paper,margin=2.3cm]{geometry}
\usepackage{amsmath,amssymb,amsthm}
\usepackage{mathtools}
\usepackage{bm}
\usepackage{booktabs}
\usepackage{array}
\usepackage{longtable}
\usepackage{enumitem}
\usepackage{xcolor}
\usepackage[most]{tcolorbox}
\usepackage{hyperref}
\usepackage{graphicx}
\usepackage{tikz}
\usepackage{microtype}
\usepackage{parskip}
\usepackage{siunitx}

\sisetup{
  output-decimal-marker={,},
  group-separator={\,},
  detect-all,
}

\hypersetup{
  colorlinks=true,
  linkcolor=blue!60!black,
  urlcolor=blue!60!black,
  citecolor=black,
  pdftitle={Metricas de calidad de mallas Q4/Q9 normalizadas a [0,1]},
  pdfauthor={EduFEM - Investigacion para reformulacion del Modulo 0},
}

\tcbuselibrary{skins,breakable}

\newtcolorbox{notabox}[1][]{
  colback=blue!4, colframe=blue!50!black, boxrule=0.6pt,
  arc=2pt, left=8pt, right=8pt, top=6pt, bottom=6pt,
  fonttitle=\bfseries, breakable, #1
}
\newtcolorbox{recbox}[1][]{
  colback=green!5, colframe=green!45!black, boxrule=0.6pt,
  arc=2pt, left=8pt, right=8pt, top=6pt, bottom=6pt,
  fonttitle=\bfseries, breakable, #1
}
\newtcolorbox{warnbox}[1][]{
  colback=orange!6, colframe=orange!60!black, boxrule=0.6pt,
  arc=2pt, left=8pt, right=8pt, top=6pt, bottom=6pt,
  fonttitle=\bfseries, breakable, #1
}

\title{\textbf{M\'etricas de calidad para cuadril\'ateros 2D}\\[2pt]
\large Normalizaci\'on a $[0,1]$, alternativas pedag\'ogicas\\ y ejemplo aplicado paso a paso}
\author{EduFEM --- Anexo de investigaci\'on para el M\'odulo 0}
\date{Mayo 2026}

\begin{document}
\maketitle

\begin{abstract}
\noindent
Este documento revisa las m\'etricas de calidad geom\'etrica m\'as utilizadas en la literatura para
cuadril\'ateros bilineales y biquadr\'aticos, propone una \textbf{normalizaci\'on uniforme a $[0,1]$}
con $1$ = elemento ideal, y desarrolla un ejemplo aplicado paso a paso sobre un Q4 trapezoidal.
Se incluye el \textit{Scaled Jacobian} junto con tres alternativas accesibles al alumno: \textit{relaci\'on
de aspecto por aristas}, \textit{esviaje angular} y \textit{taper de Robinson}. Se cierra con un
an\'alisis cruzado de qu\'e patolog\'ia geom\'etrica cubre cada combinaci\'on de m\'etricas,
proponiendo dos paquetes finales: \textbf{$\{Q_{SJ},Q_{AR}\}$} (paquete m\'inimo de 2 m\'etricas) y
\textbf{$\{Q_{SJ},Q_{AR},Q_{sk}\}$} (paquete completo de 3 m\'etricas). Las referencias siguen el
estilo Vancouver.
\end{abstract}

\tableofcontents
\vspace{6pt}\hrule\vspace{10pt}

% ====================================================================
\section{Por qu\'e normalizar a $[0,1]$}
% ====================================================================

Las m\'etricas geom\'etricas \textit{crudas} (Jacobiano determinante, relaci\'on de aspecto, \'angulos
internos, taper de Robinson) viven en dominios heterog\'eneos: el Jacobiano determinante crece con
la escala del elemento, la relaci\'on de aspecto vive en $[1,\infty)$, los \'angulos en $[0\si{\degree},180\si{\degree}]$.
Compararlas, combinarlas o presentarlas al alumno requiere homogeneizar el rango.

\subsection{Rangos candidatos}

\begin{longtable}{p{2.6cm} p{4.2cm} p{8.5cm}}
\toprule
\textbf{Rango} & \textbf{Convenci\'on} & \textbf{Discusi\'on} \\
\midrule
\endhead
$[-1,1]$ & $-1$ invertido, $0$ degenerado, $1$ ideal & Es el dominio natural del \textit{Scaled
Jacobian} (sin$\theta$). Permite distinguir invertido vs degenerado pero introduce un cero
\emph{interior} que dificulta la combinaci\'on multiplicativa y el mapeo crom\'atico mon\'otono.\\
$[0,\infty)$ & $0$ ideal, $\infty$ degenerado & Dominio natural de AR y taper. Asimetr\'ico: un
mismo orden de magnitud (4, 40, 400) tiene significados muy distintos. Mal adaptado a barras de
escala finitas.\\
$[0,100]$ & Porcentaje, $100\%$ ideal & Usado en Cubit/Trelis para reportes \cite{Stimpson2007}.
Pedag\'ogicamente equivalente a $[0,1]$ con un factor de $100$; introduce ruido visual en
gr\'aficos.\\
\rowcolor{green!7}\textbf{$\bm{[0,1]}$} & \textbf{$\bm{0}$ p\'esimo, $\bm{1}$ ideal} & \textbf{Convenci\'on
recomendada}: lectura directa, combinable como producto o m\'inimo, mapeo crom\'atico uniforme,
estandarizada por la librer\'ia \textit{Verdict} de Sandia \cite{Stimpson2007} y adoptada por
ANSYS \cite{ANSYS2021}, Abaqus, Gmsh y la mayor\'ia de los textos de generaci\'on de mallas
\cite{Knupp2001,Field2000}.\\
\bottomrule
\end{longtable}

\subsection{Justificaci\'on detallada de $[0,1]$ con $1$ = ideal}

\begin{recbox}[title=Por qu\'e $[0,1]$ es la elecci\'on \'optima]
\begin{enumerate}[leftmargin=14pt,itemsep=2pt]
  \item \textbf{Universalidad}. \textit{Verdict} (Sandia) \cite{Stimpson2007},
  ANSYS Meshing \cite{ANSYS2021} y Cubit reportan calidad en $[0,1]$ con $1$ = ideal. El alumno
  encontrar\'a la misma convenci\'on en cualquier suite profesional.
  \item \textbf{Combinabilidad}. Para reducir varias m\'etricas a un \'unico indicador de salud:
  \[
    Q_{\text{global}} = \min(Q_1,\,Q_2,\,\dots) \qquad\text{o}\qquad
    Q_{\text{prod}} = \prod_i Q_i.
  \]
  Ambas operaciones preservan $1$ como ideal y $0$ como peor caso.
  \item \textbf{Mapeo crom\'atico mon\'otono}. Una sola interpolaci\'on rojo$\to$amarillo$\to$verde
  cubre el espectro completo. Sin $0$ interior se evita la discontinuidad visual entre ``elemento
  invertido'' y ``elemento ideal por el lado equivocado''.
  \item \textbf{Inversi\'on como bandera dura}. Cuando una m\'etrica cruda da valor negativo
  (det$J \le 0$, $\sin\theta < 0$) se aplica \textit{clamp} a $0$ y se levanta una \textit{bandera
  l\'ogica} aparte (\texttt{is\_inverted=True}). Esto separa ``calidad continua'' (qu\'e tan bueno
  es un elemento sano) de ``validez bin\'aria'' (es integrable o no).
  \item \textbf{Pedagog\'ico}. La frase ``el elemento tiene calidad $0{,}78$'' se interpreta sin
  necesidad de explicar el rango.
\end{enumerate}
\end{recbox}

\begin{warnbox}[title=Trade-off conocido]
Normalizar pierde la distinci\'on \emph{cualitativa} entre ``degenerado'' (det$J=0$) y
``invertido'' (det$J<0$). Ambos quedan en $Q=0$. Por eso se recomienda complementar con una
bandera $\texttt{is\_inverted}$ que el m\'odulo educativo muestre con un icono distinto (en
EduFEM ya existe el anillo rojo discontinuo de M2). El alumno aprende que ``$Q=0$'' significa
``inservible'' sin perder la matiz f\'isica.
\end{warnbox}

% ====================================================================
\section{M\'etricas evaluadas y su normalizaci\'on}
% ====================================================================

A continuaci\'on se presentan cuatro m\'etricas representativas con su definici\'on cruda y la
versi\'on normalizada a $[0,1]$. Los s\'imbolos refieren a un cuadril\'atero $\Omega_e$ de
v\'ertices ordenados en sentido antihorario $\{N_1,N_2,N_3,N_4\}$ con coordenadas
$\bm{r}_i = (x_i,y_i)$.

\subsection{Scaled Jacobian, $Q_{SJ}$}

\textbf{Definici\'on cruda} (Knupp \cite{Knupp2001,Knupp2000a}). En cada v\'ertice $i$ se forman
los dos vectores de arista emanantes:
\[
  \bm{a}_i = \bm{r}_{i+1}-\bm{r}_i, \qquad \bm{b}_i = \bm{r}_{i-1}-\bm{r}_i.
\]
El Jacobiano local en el v\'ertice es $\bm{J}_i = [\bm{a}_i \mid \bm{b}_i]$ y se define
\[
  \mathrm{SJ}_i \;=\; \frac{\det \bm{J}_i}{\lVert\bm{a}_i\rVert\,\lVert\bm{b}_i\rVert}
  \;=\; \sin\theta_i \quad\text{(cuadril\'atero convexo CCW)},
\]
\[
  \mathrm{SJ} \;=\; \min_{i=1,\dots,4}\mathrm{SJ}_i \quad\in\;[-1,1].
\]

\textbf{Normalizaci\'on a $[0,1]$}:
\begin{equation}\boxed{\;
  Q_{SJ} \;=\; \max\!\bigl(0,\;\mathrm{SJ}\bigr) \;\in\;[0,1].
\;}\label{eq:QSJ}\end{equation}

\textbf{Concepto que el alumno reconoce}: \emph{\textquotedblleft sin\,(\'angulo interno
m\'as desviado)\textquotedblright}. En un cuadrado los cuatro \'angulos valen \SI{90}{\degree} y
$Q_{SJ}=\sin\SI{90}{\degree}=1$. Cuando un \'angulo se acerca a \SI{0}{\degree} o
\SI{180}{\degree}, su seno cae a cero.

\textbf{Cobertura}: distorsi\'on angular, casi-degeneraci\'on, inversi\'on (detectada como signo).

\medskip

\textbf{Variante por puntos de Gauss} (la usada en EduFEM): se eval\'ua
$\mathrm{SJ}_g = \det\bm{J}(\xi_g,\eta_g) / (\lVert\bm{J}_{:,0}\rVert\,\lVert\bm{J}_{:,1}\rVert)$
en los $n_g$ puntos de Gauss del elemento y se toma el m\'inimo. Para un Q4 con cuadratura $2\times2$,
el promedio entre el SJ por v\'ertice y el SJ por Gauss difiere t\'ipicamente en menos del 5\%.
La versi\'on por Gauss es m\'as conservadora porque eval\'ua \emph{donde se va a integrar la rigidez}
\cite{Knupp2000b}.

% ----------------------------------------------------------------
\subsection{Relaci\'on de aspecto, $Q_{AR}$}

\textbf{Definici\'on cruda} (Cook et al. \cite{Cook2002}). Sea $L_i = \lVert\bm{r}_{i+1}-\bm{r}_i\rVert$
la longitud de la arista $i$:
\[
  \mathrm{AR} \;=\; \frac{\max_i L_i}{\min_i L_i} \;\in\;[1,\infty).
\]

\textbf{Normalizaci\'on a $[0,1]$} (\textit{rec\'iproco simple}):
\begin{equation}\boxed{\;
  Q_{AR} \;=\; \frac{\min_i L_i}{\max_i L_i} \;\in\;(0,1].
\;}\label{eq:QAR}\end{equation}

\textbf{Concepto que el alumno reconoce}: \emph{\textquotedblleft cu\'anto m\'as largo es el lado
m\'as largo respecto al m\'as corto\textquotedblright}. Es geom\'etricamente trivial: medir cuatro
lados con regla y dividir.

\textbf{Cobertura}: elongaci\'on (psima condici\'on num\'erica de $\bm{K}_e$, anisotrop\'ia de
respuesta).

\medskip

\textbf{Variante de Robinson} \cite{Robinson1987}: en lugar de usar las aristas directamente, se
construyen los vectores $\bm{X}_1,\bm{X}_2$ del modo bilineal (definidos en
Secci\'on \ref{sec:ejemplo}). Da un valor algo distinto pero la lectura es la misma. Para uso
educativo basta con la versi\'on de aristas; la de Robinson aparece como complemento.

% ----------------------------------------------------------------
\subsection{Esviaje angular, $Q_{sk}$}

\textbf{Definici\'on cruda} (Knupp \cite{Knupp2001}, ANSYS \cite{ANSYS2021}). Para los cuatro
\'angulos internos $\theta_i$:
\[
  S_{\max} \;=\; \max_i \bigl|\theta_i - 90\si{\degree}\bigr| \;\in\;[0\si{\degree},90\si{\degree}].
\]

\textbf{Normalizaci\'on a $[0,1]$}:
\begin{equation}\boxed{\;
  Q_{sk} \;=\; 1 - \frac{S_{\max}}{90\si{\degree}} \;\in\;[0,1].
\;}\label{eq:Qsk}\end{equation}

\textbf{Concepto que el alumno reconoce}: \emph{\textquotedblleft cu\'an cerca de \SI{90}{\degree}
est\'a el \'angulo m\'as raro\textquotedblright}. Concepto netamente geom\'etrico, accesible desde
trigonometr\'ia b\'asica.

\textbf{Cobertura}: distorsi\'on angular pura. Distinta de $Q_{SJ}$ en que penaliza linealmente la
desviaci\'on respecto a \SI{90}{\degree} (el seno la penaliza cuadr\'aticamente cerca del
\'optimo).

% ----------------------------------------------------------------
\subsection{Taper de Robinson, $Q_T$}

\textbf{Definici\'on cruda} (Robinson \cite{Robinson1987}). Sobre los vectores
\[
  \bm{X}_1 = \tfrac{1}{4}(-\bm{r}_1+\bm{r}_2+\bm{r}_3-\bm{r}_4),\quad
  \bm{X}_2 = \tfrac{1}{4}(-\bm{r}_1-\bm{r}_2+\bm{r}_3+\bm{r}_4),\quad
  \bm{X}_3 = \tfrac{1}{4}( \bm{r}_1-\bm{r}_2+\bm{r}_3-\bm{r}_4),
\]
se define el \textit{taper}
\[
  T \;=\; \max\!\left( \frac{|\bm{X}_3\!\cdot\!\bm{X}_1|}{\lVert\bm{X}_1\rVert^2},\;
                       \frac{|\bm{X}_3\!\cdot\!\bm{X}_2|}{\lVert\bm{X}_2\rVert^2} \right) \;\in\;[0,1].
\]
$T=0$ corresponde a un paralelogramo; $T\!\to\!1$ a un elemento muy trapezoidal.

\textbf{Normalizaci\'on a $[0,1]$}:
\begin{equation}\boxed{\;
  Q_T \;=\; 1 - \min(1,T) \;\in\;[0,1].
\;}\label{eq:QT}\end{equation}

\textbf{Concepto que el alumno reconoce}: \emph{\textquotedblleft cu\'anto se parece a un
paralelogramo\textquotedblright}. Es la m\'etrica m\'as t\'ecnica de las cuatro porque requiere
introducir los vectores $\bm{X}_i$ de la descomposici\'on de Robinson, que no son evidentes a
primera vista.

\textbf{Cobertura}: trapezoidalidad espec\'ifica (\textit{trapezoidal locking} en Q4
\cite{MacNeal1989}).

% ====================================================================
\section{Ejemplo aplicado paso a paso}\label{sec:ejemplo}
% ====================================================================

Se considera un cuadril\'atero Q4 \emph{convexo trapezoidal} con coordenadas
\begin{center}
\begin{tabular}{cccc}
\toprule
$N_1$ & $N_2$ & $N_3$ & $N_4$ \\
\midrule
$(0,0)$ & $(4,0)$ & $(5,3)$ & $(1,4)$ \\
\bottomrule
\end{tabular}
\end{center}
ordenadas en sentido antihorario.

\begin{figure}[h]
\centering
\begin{tikzpicture}[scale=0.85]
  % grid
  \draw[gray!25, very thin] (-0.5,-0.5) grid (6,5);
  \draw[->,gray!60] (-0.5,0) -- (6,0) node[right]{$x$};
  \draw[->,gray!60] (0,-0.5) -- (0,5) node[above]{$y$};
  % nodes
  \coordinate (N1) at (0,0);
  \coordinate (N2) at (4,0);
  \coordinate (N3) at (5,3);
  \coordinate (N4) at (1,4);
  % polygon
  \fill[blue!8] (N1) -- (N2) -- (N3) -- (N4) -- cycle;
  \draw[blue!60!black, thick] (N1) -- (N2) -- (N3) -- (N4) -- cycle;
  % node markers + labels
  \foreach \p/\lab/\pos in {N1/N_1/{below left}, N2/N_2/{below right},
                            N3/N_3/{above right}, N4/N_4/{above left}}{
    \fill[blue!60!black] (\p) circle (2.4pt);
    \node[\pos, font=\small] at (\p) {$\lab$};
  }
  % edge labels (lengths)
  \node[font=\scriptsize, below=2pt] at (2,0) {$L_{12}=4$};
  \node[font=\scriptsize, right=2pt] at (4.5,1.5) {$L_{23}=\sqrt{10}$};
  \node[font=\scriptsize, above=2pt] at (3,3.5) {$L_{34}=\sqrt{17}$};
  \node[font=\scriptsize, left=2pt]  at (0.5,2) {$L_{41}=\sqrt{17}$};
\end{tikzpicture}
\caption{Cuadril\'atero Q4 del ejemplo. \'Area = $14{,}5$, todos los \'angulos internos en
$[\SI{75}{\degree},\SI{109}{\degree}]$, dos aristas iguales y dos distintas: combinaci\'on
fertil para ejercitar las cuatro m\'etricas.}
\end{figure}

% --- 3.1 -----
\subsection{Paso 1: longitudes de arista}

\[
  L_{12} \!=\! 4{,}0000, \qquad
  L_{23} \!=\! \sqrt{10} \!\approx\! 3{,}1623, \qquad
  L_{34} \!=\! \sqrt{17} \!\approx\! 4{,}1231, \qquad
  L_{41} \!=\! \sqrt{17} \!\approx\! 4{,}1231.
\]
\[
  L_{\min} = L_{23} = \sqrt{10}, \qquad L_{\max} = L_{34} = L_{41} = \sqrt{17}.
\]

% --- 3.2 -----
\subsection{Paso 2: \'angulos internos}

Para cada v\'ertice, $\cos\theta_i = (\bm{v}_{\text{prev}}\!\cdot\!\bm{v}_{\text{next}}) /
(\lVert\bm{v}_{\text{prev}}\rVert\lVert\bm{v}_{\text{next}}\rVert)$.

\medskip

\noindent\textbf{V\'ertice $N_1$:} $\bm{v}_{\text{prev}}=N_4-N_1=(1,4)$,
$\bm{v}_{\text{next}}=N_2-N_1=(4,0)$.
\[
  \cos\theta_1 = \frac{(1)(4)+(4)(0)}{\sqrt{17}\cdot 4} = \frac{1}{\sqrt{17}}
  \;\Rightarrow\; \theta_1 \approx \SI{75,964}{\degree},\; \sin\theta_1 \approx 0{,}9701.
\]

\noindent\textbf{V\'ertice $N_2$:} $\bm{v}_{\text{prev}}=N_1-N_2=(-4,0)$,
$\bm{v}_{\text{next}}=N_3-N_2=(1,3)$.
\[
  \cos\theta_2 = \frac{(-4)(1)+(0)(3)}{4\sqrt{10}} = -\frac{1}{\sqrt{10}}
  \;\Rightarrow\; \theta_2 \approx \SI{108,435}{\degree},\; \sin\theta_2 \approx 0{,}9487.
\]

\noindent\textbf{V\'ertice $N_3$:} $\bm{v}_{\text{prev}}=N_2-N_3=(-1,-3)$,
$\bm{v}_{\text{next}}=N_4-N_3=(-4,1)$.
\[
  \cos\theta_3 = \frac{(-1)(-4)+(-3)(1)}{\sqrt{10}\sqrt{17}} = \frac{1}{\sqrt{170}}
  \;\Rightarrow\; \theta_3 \approx \SI{85,601}{\degree},\; \sin\theta_3 \approx 0{,}9971.
\]

\noindent\textbf{V\'ertice $N_4$:} $\bm{v}_{\text{prev}}=N_3-N_4=(4,-1)$,
$\bm{v}_{\text{next}}=N_1-N_4=(-1,-4)$.
\[
  \cos\theta_4 = \frac{(4)(-1)+(-1)(-4)}{\sqrt{17}\sqrt{17}} = 0
  \;\Rightarrow\; \theta_4 = \SI{90,000}{\degree},\; \sin\theta_4 = 1{,}0000.
\]

\noindent\textbf{Verificaci\'on}: $\sum\theta_i = 75{,}964 + 108{,}435 + 85{,}601 + 90{,}000 = \SI{360,000}{\degree}$ \checkmark.

% --- 3.3 -----
\subsection{Paso 3: $Q_{SJ}$ por v\'ertices}

Como el cuadril\'atero es convexo y orientado CCW, $\mathrm{SJ}_i = \sin\theta_i$ y:
\[
  \mathrm{SJ}_1 = 0{,}9701, \quad
  \mathrm{SJ}_2 = 0{,}9487, \quad
  \mathrm{SJ}_3 = 0{,}9971, \quad
  \mathrm{SJ}_4 = 1{,}0000.
\]
\[
  \mathrm{SJ} = \min_i \mathrm{SJ}_i = \mathrm{SJ}_2 = 0{,}9487.
\]
Aplicando \eqref{eq:QSJ}:
\[
  \boxed{\;Q_{SJ} = \max(0,\,0{,}9487) = 0{,}9487.\;}
\]
\textbf{Lectura}: el elemento est\'a a un $5\%$ del cuadrado angular perfecto. La distorsi\'on
proviene de la obtusidad en $N_2$.

\medskip

\begin{notabox}[title=Nota: $Q_{SJ}$ por v\'ertices vs por Gauss]
La f\'ormula por v\'ertices da $Q_{SJ}=0{,}9487$. Si se eval\'ua en los $2\times 2$ puntos de
Gauss interiores (versi\'on usada por EduFEM, ec. \cite{Knupp2000b}), el resultado es
$Q_{SJ}^{(g)} = 0{,}9821$. La diferencia se debe a que los Gauss interiores ``ven'' el promedio
del Jacobiano y no el peor v\'ertice. Ambos valores son aceptables como m\'etrica; la versi\'on
por v\'ertices es m\'as estricta y m\'as f\'acil de derivar a mano, por lo que se prefiere para
el ejemplo did\'actico.
\end{notabox}

% --- 3.4 -----
\subsection{Paso 4: $Q_{AR}$ por aristas}

Aplicando \eqref{eq:QAR}:
\[
  Q_{AR} = \frac{L_{\min}}{L_{\max}} = \frac{\sqrt{10}}{\sqrt{17}}
         = \sqrt{\frac{10}{17}} \approx 0{,}7670.
\]
\[
  \boxed{\;Q_{AR} = 0{,}7670.\;}
\]
\textbf{Lectura}: el lado m\'as corto mide aproximadamente el $77\%$ del lado m\'as largo.

% --- 3.5 -----
\subsection{Paso 5: $Q_{sk}$ esviaje angular}

Desviaciones respecto a \SI{90}{\degree}:
\[
  |\theta_1-90| \!=\! 14{,}036,\quad
  |\theta_2-90| \!=\! 18{,}435,\quad
  |\theta_3-90| \!=\!  4{,}399,\quad
  |\theta_4-90| \!=\!  0{,}000\;\;(\text{en grados}).
\]
$S_{\max}=\SI{18,435}{\degree}$. Aplicando \eqref{eq:Qsk}:
\[
  Q_{sk} = 1 - \frac{18{,}435}{90} = 1 - 0{,}2048 = 0{,}7952.
\]
\[
  \boxed{\;Q_{sk} = 0{,}7952.\;}
\]
\textbf{Lectura}: el \'angulo peor desviado est\'a a un $20\%$ del \'angulo recto.

% --- 3.6 -----
\subsection{Paso 6: $Q_T$ taper de Robinson}

Vectores de la descomposici\'on bilineal:
\begin{align*}
  \bm{X}_1 &= \tfrac{1}{4}\bigl(-(0,0)+(4,0)+(5,3)-(1,4)\bigr) = \tfrac{1}{4}(8,-1) = (2{,}000,\,-0{,}250),\\
  \bm{X}_2 &= \tfrac{1}{4}\bigl(-(0,0)-(4,0)+(5,3)+(1,4)\bigr) = \tfrac{1}{4}(2,7) = (0{,}500,\,1{,}750),\\
  \bm{X}_3 &= \tfrac{1}{4}\bigl((0,0)-(4,0)+(5,3)-(1,4)\bigr) = \tfrac{1}{4}(0,-1) = (0{,}000,\,-0{,}250).
\end{align*}
Normas cuadr\'aticas:
\[
  \lVert\bm{X}_1\rVert^2 = 4{,}0625, \qquad \lVert\bm{X}_2\rVert^2 = 3{,}3125.
\]
Productos escalares con $\bm{X}_3$:
\[
  \bm{X}_3\!\cdot\!\bm{X}_1 = (0)(2)+(-0{,}25)(-0{,}25) = 0{,}0625,
\]
\[
  \bm{X}_3\!\cdot\!\bm{X}_2 = (0)(0{,}5)+(-0{,}25)(1{,}75) = -0{,}4375.
\]
\[
  t_1 = \frac{|0{,}0625|}{4{,}0625} = 0{,}0154, \qquad
  t_2 = \frac{|-0{,}4375|}{3{,}3125} = 0{,}1321.
\]
\[
  T = \max(t_1,t_2) = 0{,}1321 \;\Rightarrow\; Q_T = 1 - 0{,}1321 = 0{,}8679.
\]
\[
  \boxed{\;Q_T = 0{,}8679.\;}
\]
\textbf{Lectura}: el elemento es \emph{casi} un paralelogramo (un paralelogramo perfecto dar\'ia
$Q_T=1$).

% --- 3.7 -----
\subsection{Paso 7: tabla resumen y diagn\'ostico}

\begin{center}
\renewcommand{\arraystretch}{1.15}
\begin{tabular}{l c c l}
\toprule
\textbf{M\'etrica} & \textbf{Valor cruda} & \textbf{$Q\in[0,1]$} & \textbf{Lectura} \\
\midrule
Scaled Jacobian (m\'in v\'ert.)  & $0{,}9487$           & $0{,}9487$ & calidad angular muy buena \\
SJ por Gauss (proy. EduFEM)      & $0{,}9821$           & $0{,}9821$ & m\'as benevolente (interior) \\
Relaci\'on de aspecto            & $\mathrm{AR}=1{,}3038$ & $0{,}7670$ & lados con $30\%$ de variaci\'on \\
Esviaje angular                  & $S_{\max}=\SI{18,4}{\degree}$ & $0{,}7952$ & \'angulos a $20\%$ del recto \\
Taper de Robinson                & $T=0{,}1321$         & $0{,}8679$ & casi paralelogramo \\
\bottomrule
\end{tabular}
\end{center}

\noindent\textbf{Indicador combinado m\'inimo} (m\'as conservador):
\[
  Q_{\min} = \min(Q_{SJ},Q_{AR},Q_{sk},Q_T) = Q_{AR} = 0{,}7670.
\]
\textbf{Indicador producto} (m\'as severo):
\[
  Q_{\prod} = 0{,}9487 \cdot 0{,}7670 \cdot 0{,}7952 \cdot 0{,}8679 \approx 0{,}5023.
\]

\textbf{Diagn\'ostico global}: el elemento es \emph{aceptable} (Q$_{\min}\approx 0{,}77$, fuera del
rango cr\'itico $<\!0{,}5$). El factor m\'as limitante es la \emph{relaci\'on de aspecto}: la
arista $N_2 N_3$ es $30\%$ m\'as corta que las laterales. La calidad angular y la
trapezoidalidad son muy buenas.

% ====================================================================
\section{Combinaciones y cobertura de patolog\'ias}
% ====================================================================

Existen \emph{cinco patolog\'ias} geom\'etricas que pueden degradar el desempe\~no de un Q4 o Q9:

\begin{enumerate}[leftmargin=14pt,itemsep=2pt]
  \item \textbf{Inversi\'on} (det$J \le 0$): la matriz $\bm{B}$ no se puede calcular; integraci\'on
  imposible.
  \item \textbf{Distorsi\'on angular} (\'angulos muy lejos de \SI{90}{\degree}): degrada el
  condicionamiento y la precisi\'on de tensiones.
  \item \textbf{Elongaci\'on} (relaci\'on de aspecto alta): empeora $\kappa(\bm{K})$; introduce
  anisotrop\'ia de respuesta.
  \item \textbf{Trapezoidalidad} (lados opuestos no paralelos): activa el \textit{trapezoidal
  locking} en Q4 \cite{MacNeal1989,Cook2002}.
  \item \textbf{Curvatura interna Q9} (mid-nodes fuera del cuarto medio o centroide fuera de la
  envolvente): invalida el mapeo isoparam\'etrico Q9. \emph{Fuera del alcance de este documento}
  porque el M\'odulo 0 reformulado opera sobre los cuatro v\'ertices macro y, como aclara el
  enunciado, no manipula los nodos medios Q9.
\end{enumerate}

\subsection{Matriz cruzada m\'etrica $\times$ patolog\'ia}

\begin{center}
\renewcommand{\arraystretch}{1.2}
\begin{tabular}{l c c c c}
\toprule
\textbf{M\'etrica} & \textbf{Inv.} & \textbf{Distor.\ ang.} & \textbf{Elongaci\'on} & \textbf{Trapezoid.} \\
\midrule
$Q_{SJ}$  & \checkmark fuerte & \checkmark fuerte & --- & parcial \\
$Q_{AR}$  & ---               & ---                & \checkmark fuerte & --- \\
$Q_{sk}$  & parcial           & \checkmark fuerte & ---               & parcial \\
$Q_T$     & ---               & ---                & ---               & \checkmark fuerte \\
\bottomrule
\end{tabular}
\end{center}

\subsection{Cobertura por paquete}

\begin{longtable}{p{4cm} p{11cm}}
\toprule
\textbf{Paquete} & \textbf{Cobertura y observaciones} \\
\midrule
$\{Q_{SJ}\}$ &
Detecta inversi\'on y distorsi\'on angular fuerte. \textbf{No detecta} un rect\'angulo
$1\!\times\!100$ (los \'angulos siguen siendo \SI{90}{\degree}, $Q_{SJ}=1$) ni un trapecio con
\'angulos rectos. Insuficiente como m\'etrica \'unica.\\[2pt]
\rowcolor{green!7}\textbf{$\{Q_{SJ},Q_{AR}\}$} \textsc{(m\'inimo)} &
Cubre las tres patolog\'ias m\'as comunes: \textbf{inversi\'on, distorsi\'on angular y
elongaci\'on}. Deja sin tratar trapezoidalidad pura, pero \'esta es un caso poco frecuente en
mallas estructuradas educativas. \textbf{Recomendado como paquete pedag\'ogico m\'inimo}.\\[2pt]
\rowcolor{green!12}\textbf{$\{Q_{SJ},Q_{AR},Q_{sk}\}$} \textsc{(completo)} &
A\~nade un indicador angular \emph{lineal} (mientras $Q_{SJ}$ es cuadr\'atico cerca del
\'optimo). Cuando el alumno desliza un nodo, $Q_{sk}$ reacciona m\'as r\'apido que $Q_{SJ}$;
pedag\'ogicamente m\'as informativo. \textbf{Recomendado como paquete completo}. Sigue sin
detectar trapezoidalidad pura, pero cubre $\ge 95\%$ de los casos reales.\\[2pt]
$\{Q_{SJ},Q_{AR},Q_{sk},Q_T\}$ \textsc{(exhaustivo)} &
Cubre las cuatro patolog\'ias 2D. Incluir $Q_T$ obliga a introducir los vectores de Robinson:
poco intuitivo. Recomendado solo para una variante avanzada del m\'odulo o para un m\'odulo
posterior sobre \textit{locking}.\\
\bottomrule
\end{longtable}

\begin{recbox}[title=Conclusi\'on operativa]
\textbf{Para EduFEM (M\'odulo 0 reformulado)} se recomienda el paquete
$\{Q_{SJ}, Q_{AR}\}$ como configuraci\'on por defecto, con la opci\'on de habilitar $Q_{sk}$
como tercera m\'etrica para alumnos que quieran un diagn\'ostico m\'as fino. Esta combinaci\'on:
\begin{itemize}[leftmargin=14pt,itemsep=1pt]
  \item Usa \'unicamente conceptos de trigonometr\'ia b\'asica (seno, longitud de segmento).
  \item Cubre las tres patolog\'ias geom\'etricas m\'as relevantes en pr\'actica.
  \item Permite una barra de calidad horizontal por m\'etrica con escala embebida.
  \item Mantiene la representaci\'on crom\'atica del M\'odulo 0 (verde/amarillo/rojo).
  \item Aplica indistintamente a Q4 y Q9 sin diferenciar tipos (lo cual era expl\'icitamente
  parte del requerimiento).
\end{itemize}
\end{recbox}

% ====================================================================
\section{Conclusiones}
% ====================================================================

\begin{enumerate}[leftmargin=14pt]
  \item \textbf{Rango}: el dominio $[0,1]$ con $1$ = ideal es la opci\'on que mejor balancea
  legibilidad, combinabilidad y compatibilidad con software comercial. Se justifica por la
  convenci\'on Verdict/ANSYS \cite{Stimpson2007,ANSYS2021}, por la posibilidad de combinar
  m\'etricas por m\'inimo o producto, y por el mapeo crom\'atico mon\'otono que ya emplea el
  M\'odulo 0.
  \item \textbf{Inversi\'on como bandera externa}: el \textit{clamp} a $0$ se complementa con un
  flag \texttt{is\_inverted} que activa una alarma distinta (anillo rojo discontinuo, como ya
  hace M2). De este modo $Q\!=\!0$ significa siempre ``inservible'' sin perder la matiz f\'isica
  entre degenerado e invertido.
  \item \textbf{Paquete recomendado}: $\{Q_{SJ}, Q_{AR}\}$ como m\'inimo did\'actico,
  $\{Q_{SJ}, Q_{AR}, Q_{sk}\}$ como completo. Ambas opciones usan conceptos accesibles
  (seno de un \'angulo, longitud de un segmento, desviaci\'on respecto a \SI{90}{\degree}) y son
  geom\'etricamente verificables a mano.
  \item \textbf{Lo que el paquete NO cubre}: la trapezoidalidad pura (un Q4 con \'angulos rectos
  pero lados opuestos no paralelos) y la admisibilidad de mid-nodes Q9. La primera puede
  abordarse en un m\'odulo posterior dedicado a \textit{shear/trapezoidal locking}; la segunda
  est\'a fuera del alcance del M\'odulo 0 reformulado seg\'un el requerimiento del usuario.
\end{enumerate}

% ====================================================================
\section*{Referencias} % Vancouver (numeric, by first appearance)
% ====================================================================

\addcontentsline{toc}{section}{Referencias}
\renewcommand{\arraystretch}{1.05}
\begin{enumerate}[label={[\arabic*]},leftmargin=24pt,itemsep=2pt]

\item\label{Stimpson2007} Stimpson CJ, Ernst CD, Knupp PM, P\'ebay PP, Thompson D. The Verdict
geometric quality library. Albuquerque (NM): Sandia National Laboratories; 2007. Report No.\
SAND2007-1751.

\item\label{ANSYS2021} ANSYS Inc. ANSYS Meshing User's Guide, Release 2021 R1. Canonsburg (PA):
ANSYS Inc.; 2021. Cap\'itulo 17: Mesh Quality Metrics.

\item\label{Knupp2001} Knupp PM. Algebraic mesh quality metrics. SIAM Journal on Scientific
Computing. 2001;23(1):193--218. doi:10.1137/S1064827500371499.

\item\label{Knupp2000a} Knupp PM. Achieving finite element mesh quality via optimization of the
Jacobian matrix norm and associated quantities. Part I --- a framework for surface mesh
optimization. International Journal for Numerical Methods in Engineering.
2000;48(3):401--420.

\item\label{Knupp2000b} Knupp PM. Achieving finite element mesh quality via optimization of the
Jacobian matrix norm and associated quantities. Part II --- a framework for volume mesh
optimization. International Journal for Numerical Methods in Engineering.
2000;48(8):1165--1185.

\item\label{Field2000} Field DA. Qualitative measures for initial meshes. International Journal
for Numerical Methods in Engineering. 2000;47(4):887--906.

\item\label{Robinson1987} Robinson J. CRE method of element testing and the Jacobian shape
parameters. Engineering Computations. 1987;4(2):113--118.

\item\label{Cook2002} Cook RD, Malkus DS, Plesha ME, Witt RJ. Concepts and Applications of Finite
Element Analysis. 4th ed. Hoboken (NJ): John Wiley \& Sons; 2002. Cap\'itulos 6 y 7.

\item\label{MacNeal1989} MacNeal RH. Toward a defect-free four-noded membrane element.
Finite Elements in Analysis and Design. 1989;5(1):31--37.

\item\label{Bathe2014} Bathe KJ. Finite Element Procedures. 2nd ed. Watertown (MA): Klaus-J\"urgen
Bathe; 2014.

\item\label{Zienkiewicz2013} Zienkiewicz OC, Taylor RL, Zhu JZ. The Finite Element Method: Its
Basis and Fundamentals. 7th ed. Oxford: Butterworth-Heinemann; 2013.

\item\label{Hughes2000} Hughes TJR. The Finite Element Method: Linear Static and Dynamic Finite
Element Analysis. Mineola (NY): Dover Publications; 2000.

\item\label{Lee1994} Lee NS, Bathe KJ. Effects of element distortions on the performance of
isoparametric elements. International Journal for Numerical Methods in Engineering.
1994;36(20):3553--3576.

\end{enumerate}

\bigskip

\hrule
\medskip
\noindent\footnotesize\emph{Documento generado para el proyecto \textbf{EduFEM} como
fundamentaci\'on de la reformulaci\'on del M\'odulo 0 (Calidad de la malla). Las normalizaciones
y el ejemplo aplicado se computaron contra las funciones \texttt{scaled\_jacobian},
\texttt{robinson\_aspect}, \texttt{robinson\_taper} e \texttt{internal\_angles} de
\texttt{fem/mesh\_quality.py}.}

\end{document}
